% ============================================================
%  Capítulo 1 — Introdução
% ============================================================
\chapter{INTRODUÇÃO}
\label{cap:intro}

\section{Contextualização e motivação}
\label{sec:contexto}

A rastreabilidade de usuários na web deixou de depender exclusivamente de
cookies há mais de uma década. Técnicas de browser fingerprinting
(identificação de usuários por meio de atributos coletados silenciosamente pelo
navegador) tornaram-se o principal mecanismo de rastreamento persistente
utilizado por scripts de terceiros, redes de anúncios e serviços de análise de
audiência \cite{eckersley2010}. Ao contrário dos cookies, o
fingerprint não é armazenado no dispositivo do usuário e não pode ser
apagado com a limpeza do histórico ou com o modo de navegação privada.

Um fingerprint de navegador é construído a partir da combinação de dezenas
de propriedades mensuráveis: o resultado da renderização de formas geométricas na
API Canvas, o comportamento numérico do processamento de áudio via Web Audio API,
os parâmetros de capacidade gráfica reportados pela API WebGL, as dimensões e
resolução de cor da tela, as propriedades de hardware expostas pelo objeto
\texttt{navigator} (\texttt{hardwareConcurrency}, \texttt{deviceMemory},
\texttt{languages}, \texttt{platform}) e, em alguns casos, o endereço IP local
revelado pelo protocolo WebRTC. A combinação desses vetores é suficientemente
discriminativa para identificar mais de 99\% dos visitantes de um mesmo site de
forma unívoca \cite{laperdrix2016}.

Bibliotecas de fingerprinting de código aberto, como FingerprintJS e
CreepJS, são implantadas em milhares de sites da web de alta visitação. Essas
bibliotecas empacotam os mecanismos de coleta em funções JavaScript que executam
integralmente dentro do processo do navegador, tornando a detecção por filtragem
de URL inadequada: o mesmo bundle JavaScript pode conter tanto lógica
legítima de interface quanto rotinas de identificação oculta.

A abordagem mais direta para detectar fingerprinting no nível do código é
analisar os bytecodes gerados pelo motor JavaScript para cada função, uma vez que
operações como \texttt{canvas.toDataURL}, \texttt{AudioContext.createOscillator} e
\texttt{navigator.hardwareConcurrency} produzem sequências de bytecodes observáveis
sem que seja necessário executar ou instrumentar a função em nível de aplicação.
ByteDefender \cite{nikkhahbahrami2025} demonstrou que sequências de bytecodes do
motor V8, utilizado pelo Chromium, são discriminativas o suficiente para
detectar funções de fingerprinting com F1 $\approx 0{,}95$ usando um
classificador treinado sobre uma representação bag-of-ngrams dos opcodes.

Contudo, o V8 é apenas um dos dois grandes motores JavaScript em uso em
navegadores de desktop. O Firefox utiliza o SpiderMonkey, cuja arquitetura interna
(estratégias de alocação de registradores, estrutura do grafo de controle de
fluxo, hierarquia de representações intermediárias) difere substancialmente da
do V8. Um modelo treinado sobre bytecodes brutos do V8 não transfere para o
SpiderMonkey: os conjuntos de opcodes são disjuntos, e as sequências produzidas
pelo mesmo código JavaScript diferem em estrutura e distribuição entre os dois
motores. A consequência prática é que a abordagem do ByteDefender requer
retreinamento completo para cada motor-alvo, o que representa uma barreira de
engenharia significativa para qualquer detector implantado em um ecossistema de
navegadores heterogêneo.

Este trabalho endereça diretamente essa limitação. A contribuição central é o
projeto e a validação de uma \textbf{representação intermediária (RI)
cross-engine} que mapeia os opcodes do V8 (Ignition) e do SpiderMonkey para um
vocabulário compartilhado de 145 tokens, com cobertura de 99,82\% dos tokens
reais observados em um crawl de 68 milhões de funções JavaScript. Sobre
essa RI, o LightGBM com bigramas atinge F1~=~0,843 como modelo de referência da linha principal.
Uma estratégia de inferência iterativa,
que recomputa contagens de funções previstas em múltiplos limiares como features de
contexto de script, eleva o resultado implantável para \textbf{F1~=~0,904} sem
nenhum rótulo externo. O modelo conjunto com contexto iterativo, treinado em dados de ambos os motores,
atinge F1~=~0,903 no Chromium e F1~=~0,904 no Firefox a partir de um único
conjunto de pesos, sem retreinamento por motor.

O custo da portabilidade entre motores é de aproximadamente 5 pontos de F1 em
relação ao ByteDefender (F1~=~0,904 versus $\approx$0,950), refletindo a perda de
informação causada pela abstração da RI e pelo ruído distribucional da presença de
dois motores no corpus de treinamento. Parte desse custo também decorre de uma
tarefa estritamente mais difícil: este trabalho inclui funções anônimas, sete
categorias de fingerprinting (contra quatro do ByteDefender) e utiliza
exclusivamente bytecodes sem qualquer metadado de nome de função. O modelo conjunto
exige apenas um conjunto de pesos para operar corretamente nos dois principais
navegadores de desktop.

\section{Perguntas de pesquisa}
\label{sec:perguntas}

Este trabalho é orientado pelas seguintes perguntas de pesquisa:

\textbf{PQ1.} É possível projetar uma representação intermediária cross-engine
que mapeie opcodes do V8 e do SpiderMonkey para um vocabulário comum com cobertura
superior a 99\% dos tokens observados em um corpus de escala real?

\textbf{PQ2.} Um único classificador treinado sobre essa RI consegue detectar
funções de fingerprinting em ambos os motores com F1 comparável ao estado
da arte, sem retreinamento por motor?

\textbf{PQ3.} A RI é suficiente para viabilizar transferência zero-shot entre
motores, ou o treinamento conjunto também é necessário?

\textbf{PQ4.} Qual é o teto de informação da detecção baseada exclusivamente em
bytecodes cross-engine, isto é, qual o F1 máximo alcançável antes que a
expansão do espaço de features deixe de produzir ganhos mensuráveis?

%, , , , , , , , , , , , , , , , , , , --
\section{Contribuições}
\label{sec:contribuicoes}

As contribuições deste trabalho são:

\begin{enumerate}

\item \textbf{Instrumentação do SpiderMonkey com rastreamento de funções
anônimas.} Modificação do código-fonte do Firefox/SpiderMonkey para emitir, a
cada função compilada, o URL do script, o índice interno, o nome da função e a
sequência de bytecodes. Para funções anônimas, um identificador único
(\texttt{anon\_injected\_ID}) é injetado como \texttt{displayName} da função,
propagando-se para o \texttt{Error.stack} da chamada e possibilitando o
join com as chamadas de API sem perda de dados.

\item \textbf{Pipeline de coleta cross-engine.} Crawlers
Selenium/Firefox e Playwright/Chromium que visitam os 10.000 sites mais
populares da lista Tranco com os dois navegadores instrumentados, emitindo logs no
mesmo formato JSONL unificado. Uma extensão de navegador compartilhada (Manifest
V3) intercepta 32 APIs de fingerprinting em ambos os motores e normaliza
propriedades Firefox-exclusivas, de modo que as bibliotecas de
fingerprinting tomem o mesmo caminho de código nos dois navegadores.

\item \textbf{Dataset cross-engine de escala real.} 68.008.403 funções
JavaScript rotuladas provenientes de ambos os motores, com 37.644 funções de
fingerprinting (0,055\%), cobrindo sete categorias: Canvas, Áudio, WebRTC,
Fonte, WebGL, Tela e Navigator.

\item \textbf{Representação intermediária cross-engine (contribuição principal).}
Metodologia baseada em 27 micro-scripts cobrindo construções JavaScript relevantes
para fingerprinting, com execução pareada em ambos os motores e alinhamento
manual e empírico das sequências de bytecodes. O dicionário resultante contém 198
mapeamentos V8 e 236 mapeamentos SpiderMonkey para 145 tokens de RI, com cobertura
de 99,82\% no corpus real.

\item \textbf{Avaliação comparativa de classificadores.} Treinamento e avaliação
de Random Forest, LightGBM, SVM, BiLSTM e Transformer para detecção binária (is\_fp)
e multilabel (7 categorias) sobre as sequências de RI. Experimentos de
generalização cross-engine: transferência zero-shot (F1 $\approx 0{,}15$) e modelo
conjunto com contexto iterativo (F1~=~0,903 no Chromium, F1~=~0,904 no Firefox).
Análise do teto de informação via comparação
unigrams $\to$ bigrams $\to$ trigrams. Benchmark de inferência: 44~µs/função em
CPU a lote de 1.000.

\end{enumerate}

%, , , , , , , , , , , , , , , , , , , --
\section{Estrutura do documento}
\label{sec:estrutura}

O restante deste documento está organizado da seguinte forma.

O \textbf{Capítulo~\ref{cap:fundamentacao}} apresenta a fundamentação teórica e
os trabalhos relacionados: os mecanismos de browser fingerprinting, a
arquitetura dos motores V8 e SpiderMonkey, o trabalho ByteDefender e as técnicas
de representação de código por NLP (TF-IDF, word embeddings) utilizadas
neste trabalho.

O \textbf{Capítulo~\ref{cap:arquitetura}} descreve a arquitetura do sistema: a
instrumentação do V8 e do SpiderMonkey, a extensão de navegador compartilhada, os
crawlers, a normalização de propriedades Firefox e a estratégia de
join entre bytecodes e chamadas de API.

O \textbf{Capítulo~\ref{cap:dataset}} caracteriza o dataset: configuração do
crawl, estatísticas do corpus, heurísticas de rotulagem, limpeza de ruído e
estratégia de amostragem para treinamento.

O \textbf{Capítulo~\ref{cap:ir}} detalha a representação intermediária
cross-engine: motivação, metodologia dos micro-scripts, o dicionário de
equivalência e o pipeline de tradução do dataset. Descreve também o treinamento
dos embeddings FastText sobre o corpus traduzido.

O \textbf{Capítulo~\ref{cap:experimentos}} apresenta os experimentos e
resultados: configuração experimental, resultados dos classificadores binário e
multilabel, generalização cross-engine (zero-shot e modelo conjunto) e benchmark
de inferência.

O \textbf{Capítulo~\ref{cap:discussao}} discute os resultados: o teto de
informação da RI, a necessidade do treinamento conjunto, os limites do multilabel
para rótulos raros, o custo da independência de motor em comparação com o
ByteDefender e a implantabilidade prática.

O \textbf{Capítulo~\ref{cap:conclusao}} conclui o trabalho com uma síntese das
contribuições, as limitações identificadas e as direções para trabalhos futuros.

O \textbf{Apêndice~\ref{ap:ir-dict}} apresenta o dicionário de equivalência
completo da RI cross-engine (145 tokens, 198 mapeamentos V8, 236 mapeamentos
SpiderMonkey). O \textbf{Apêndice~\ref{ap:instrumentacao}} contém os detalhes
técnicos da instrumentação do V8 e do SpiderMonkey. O
\textbf{Apêndice~\ref{ap:crawler-ext}} reúne os detalhes de implementação dos
crawlers e da extensão de navegador compartilhada.
